\documentclass[twocolumn,11pt]{article}

% Packages
\usepackage[utf8]{inputenc}
\usepackage[T1]{fontenc}
\usepackage{textcomp}
\usepackage[margin=1in]{geometry}
\usepackage{graphicx}
\usepackage{amsmath}
\usepackage{amssymb}
\usepackage{hyperref}
\usepackage{xcolor}
\usepackage{booktabs}
\usepackage{caption}
\usepackage{float}
\usepackage{parskip}
\usepackage{algorithm}
\usepackage{algorithmic}
\usepackage{titlesec}
\usepackage[numbers,sort&compress]{natbib}
\usepackage{balance}

% Better column balancing
\setlength{\columnsep}{0.25in}
\flushbottom

% Paragraph spacing: reduced to 4pt
\setlength{\parskip}{4pt}
\setlength{\parindent}{0pt}

% Section spacing: reduced to 4pt
\titlespacing*{\section}{0pt}{4pt}{4pt}
\titlespacing*{\subsection}{0pt}{4pt}{4pt}
\titlespacing*{\subsubsection}{0pt}{4pt}{4pt}

% Figure and table spacing: reduced to 4pt
\setlength{\floatsep}{4pt}
\setlength{\textfloatsep}{4pt}
\setlength{\intextsep}{4pt}

% Allow more flexible float placement
\renewcommand{\topfraction}{0.9}
\renewcommand{\bottomfraction}{0.8}
\renewcommand{\textfraction}{0.07}
\renewcommand{\floatpagefraction}{0.7}
\setlength{\floatsep}{6pt plus 2pt minus 2pt}
\setlength{\textfloatsep}{8pt plus 2pt minus 2pt}
\setlength{\intextsep}{8pt plus 2pt minus 2pt}

% Title and authors
\title{\textbf{Signal Hierarchical Petri Nets Capture Emergent Nonlinear Dynamics in MAPK Cascades: Bistability, Excitability, and Perfect Adaptation}}

\author{
Eugênio Simão\textsuperscript{1,*}\\
\textsuperscript{1}Department of Computer Science\\
Universidade Federal de Santa Catarina (UFSC)\\
Araranguá, Santa Catarina, 88906-072, Brazil\\
\textsuperscript{*}Corresponding author: eugenio.simao@ufsc.br
}

\date{January 11, 2026}

\begin{document}

\maketitle

\begin{abstract}
Cells process information through signaling cascades that can remember past states, amplify weak signals, oscillate rhythmically, or adapt to persistent stimuli. The MAPK/ERK cascade exemplifies this versatility: depending on how feedback loops are wired, the same molecular components produce qualitatively different behaviors. Using computational modeling, we show that a single three-tier phosphorylation cascade (Raf→MEK→ERK) recapitulates four canonical nonlinear phenomena through feedback architecture alone. Strong positive feedback—where active ERK accelerates its own production by degrading phosphatases—creates bistability with stable ON and OFF states separated by 577-fold, enabling cellular memory. Balanced opposing feedbacks produce excitability, where the cascade ignores small stimuli but fires all-or-nothing responses to suprathreshold inputs with 500-fold amplification. Weakened negative feedback generates sustained oscillations at 20 cycles per minute, creating pulsatile ERK activity characteristic of growth factor signaling. Most remarkably, dual-pathway negative feedback achieves near-perfect adaptation: the cascade responds transiently to growth factor pulses with 3,000-fold amplification, yet returns to baseline within seconds despite sustained stimulation. This adaptive response operates at 98.8\% thermodynamic efficiency—approaching the physical limit where steady-state elevation (118×) reaches the theoretical minimum (117×) imposed by high transient amplification. Our analysis reveals fundamental trade-offs: the same cascade cannot simultaneously exhibit memory, excitability, oscillations, and adaptation. These results suggest cells may dynamically reconfigure feedback architecture to switch between computational modes, indicating that signaling networks function as reprogrammable biological circuits rather than fixed pathways.
\end{abstract}

\textbf{Keywords:} MAPK Cascade, Bistability, Excitability, Oscillations, Perfect Adaptation, Feedback Loops, Cellular Information Processing, Systems Biology

\section*{Introduction}

Cells implement computational functions—memory, amplification, oscillation, adaptation—using the same molecular cascades. The MAPK/ERK pathway (Raf→MEK→ERK) exemplifies this versatility: in oocytes it exhibits bistability (persistent ON/OFF states \citep{ferrell2009}), in neurons excitability (all-or-nothing spikes \citep{sasagawa2005}), in proliferating cells oscillations (periodic pulsing \citep{shankaran2009}), and in hormone responses adaptation (transient sensitivity despite sustained stimulation \citep{mackeigan2005, amit2007}). How do the same kinases produce qualitatively different behaviors?

Feedback architecture provides the answer: positive feedback creates bistability, negative feedback enables adaptation. But can a single cascade support all four phenomena? What are the thermodynamic efficiency limits when amplifying 1,000× while adapting to baseline?

Here we confront a modeling limitation. Classical Petri nets (PN = $(P, T, F, W, M_0)$: places, transitions, arcs, weights, marking) elegantly capture discrete reactions and spatial organization through bipartite token-flow graphs. Yet they fundamentally restrict dependencies to arc-connected molecules. A kinase sensing extracellular growth factors requires explicit arc chains through receptors and membrane transport—cumbersome for \emph{information} sensing where cells detect signals without consuming them (quorum sensing, ATP status, hormone gradients).

This becomes critical for hierarchical control. Bacteriophage lambda stress signals (UV damage → RecA activation) \emph{preempt} developmental programs (lysogeny → CI repressor) through information-theoretic override, not mass transfer. Classical Petri nets, lacking non-local signal channels, cannot naturally model such hierarchical precedence.

The theoretical innovation: extend Petri nets from 5-tuple to 13-tuple formalism by adding \emph{signal places} ($\Psi$)—places referenced in rate functions without arc connections, enabling non-local sensing. This seemingly modest change has discovery consequences: extended Petri nets predict hierarchical cascade preemption in lambda phage (RecA 43\% mutual information overrides CII 74\% MI with 2.01× priority, validated experimentally) and multi-layer decision hierarchies in bacterial sporulation \citep{simao2025thermodynamic}. The question is whether this principle generalizes to eukaryotic signaling cascades.

We demonstrate that it does. MAPK cascades, modeled through extended Petri nets with spatial organization (membrane/cytoplasm/nucleus), thermodynamic validation ($\Delta$G constraints), and hybrid simulation (stochastic + deterministic), recapitulate all four computational modes through feedback modulation alone: bistability (577× separation, $\alpha/\beta$>2), excitability (500× amplification, $\alpha/\beta$$\approx$0.1), oscillations (20/min, $\alpha/\beta$<0.2), and adaptation (96.4\% baseline return, 3,219× transient, 98.8\% thermodynamic efficiency). Hierarchical preemption manifests through signal places: growth factor and MKP levels act as global information orchestrators crossing compartment boundaries (membrane/cytoplasm/nucleus), hierarchically determining which computational mode dominates—high GF with low MKP preempts adaptation in favor of bistability, while sustained GF with high MKP enforces adaptation. This validates the extended Petri net principle across kingdoms: signals orchestrate local deterministic decisions through information-flux hierarchy.

\section*{Results}

Lambda phage extended Petri net analysis revealed hierarchical preemption: UV stress (RecA 43\% MI) overrides developmental signals (CII 74\% MI) with 2.01× priority through explicit signal consumption—a result validated experimentally (UV→71\% lytic). If universal, this principle should extend to MAPK cascades across four computational modes: bistability (memory), excitability (thresholding), oscillations (rhythm), and adaptation (change detection). Each mode poses distinct challenges: bistability requires two stable states, excitability needs all-or-nothing thresholds, oscillations demand sustained periodicity, and adaptation must respond transiently yet return to baseline despite persistent stimulation.

We modeled Raf→MEK→ERK with spatial organization (membrane/cytoplasm/nucleus) and systematic feedback modulation.

\subsection*{Bistability: Cellular Memory Through Positive Feedback}

\textbf{The challenge:} Creating bistable behavior requires engineering a system with two stable states—ON and OFF—that resist perturbations. The cascade must "remember" past stimulation: once activated, it should remain active even after growth factors disappear. This requires positive feedback strong enough to maintain self-sustaining activation, but not so strong that the OFF state becomes unstable.

We implemented positive feedback by allowing active ERK to accelerate its own production through phosphatase degradation. Specifically, ERK-PP promotes degradation of PP2A, the phosphatase that would otherwise deactivate Raf and MEK. This creates a self-reinforcing loop: more ERK → less PP2A → more active Raf/MEK → more ERK.

Our bistable ERK cascade exhibited persistent memory of transient stimulation. Starting from the LOW state (ERK-PP $\sim$1 nM), a brief pulse of growth factor (100 nM for 30 seconds) triggered full activation to the HIGH state (ERK-PP = 577 nM, representing 96\% activation). Remarkably, when the stimulus was removed, the system remained in the HIGH state rather than returning to baseline—the cascade "remembered" the stimulation event.

\begin{figure}[htbp]
\centering
\includegraphics[width=0.9\columnwidth]{mapk_basin_of_attraction.pdf}
\caption{Bistable MAPK memory with 577-fold state separation.}
\label{fig:bistability}
\end{figure}

This memory persisted indefinitely, with the HIGH state maintained at ERK-PP levels 577-fold higher than the LOW state. The basin of attraction (Figure~\ref{fig:bistability}) reveals two stable steady states separated by an unstable threshold—different stimulus histories lead to different final states, demonstrating hysteresis. The mechanism involves ERK-PP driving degradation of the phosphatase PP2A, creating a self-reinforcing positive feedback loop where active ERK promotes its own activation by reducing dephosphorylation capacity. This autocatalytic architecture enables the cascade to maintain distinct states—a form of molecular memory encoding prior exposure to growth factors.

\subsection*{Excitability: All-or-Nothing Responses to Graded Stimuli}

\textbf{The challenge:} Excitable systems must convert graded inputs into binary outputs—sharp, stereotyped responses triggered only above a threshold. Below threshold, the system ignores the stimulus; above threshold, it fires a full-amplitude spike regardless of stimulus strength. Creating this switch-like behavior requires balancing positive feedback (amplification) against negative feedback (recovery) while maintaining a sharp threshold that distinguishes signal from noise.

Our balanced ERK cascade (positive feedback balanced with MKP-mediated negative feedback) displayed classic excitable behavior. Subthreshold growth factor doses (<10 nM) produced negligible ERK activation (~1 nM), while suprathreshold doses (>10 nM) triggered all-or-nothing spikes to 505 nM ERK-PP—a 500-fold amplification independent of stimulus strength above threshold.

\begin{figure}[htbp]
\centering
\includegraphics[width=0.9\columnwidth]{mapk_excitability.pdf}
\caption{Excitable MAPK dynamics with all-or-nothing spikes.}
\label{fig:excitability}
\end{figure}

The dose-response curve (Figure~\ref{fig:excitability}) exhibited a sharp step function characteristic of excitable systems: the amplitude of the spike remained constant once the threshold was crossed, encoding stimulus presence as a binary event rather than graded intensity. After stimulus removal, ERK-PP returned to baseline within 30 seconds via MKP-mediated feedback, enabling the cascade to "reset" and respond to subsequent stimuli. This refractoriness is essential for excitable systems to transmit information through spike sequences rather than maintaining persistent activation.

\subsection*{Oscillations: Sustained Rhythmic Signaling}

\textbf{The challenge:} Creating sustained oscillations demands a mechanism that cycles continuously without damping to steady state. Oscillations require delayed negative feedback—the inhibitor must accumulate slowly enough that the activator overshoots before being suppressed, but fast enough to bring the system back down, creating periodic cycles. The MAPK cascade, with its three-tier phosphorylation architecture and fast reaction kinetics, converges rapidly to steady states, making oscillatory behavior particularly challenging.

Despite these constraints, we achieved robust oscillations in the ERK cascade through continuous MKP synthesis feedback. The system exhibited sustained periodic behavior with ERK-PP oscillating between 0-427 nM at a frequency of 20.2 oscillations per minute (period 2.98 $\pm$ 0.30 seconds). Over 180 seconds, the cascade completed 54 oscillatory cycles (Figure~\ref{fig:oscillations}) with consistent amplitude and period—demonstrating that even rapid post-translational cascades can sustain rhythmic signaling under appropriate feedback architecture.

\begin{figure}[htbp]
\centering
\includegraphics[width=0.9\columnwidth]{mapk_oscillations_timed.pdf}
\caption{Sustained oscillatory MAPK dynamics.}
\label{fig:oscillations}
\end{figure}

The oscillatory mechanism relies on ERK-PP driving its own inhibitor (MKP) synthesis with Hill coefficient kinetics, creating a delayed negative feedback loop. When ERK-PP rises, MKP synthesis accelerates, but with a delay allowing ERK-PP to peak before MKP accumulation suppresses it. As MKP degrades, ERK-PP rises again, initiating the next cycle. The pulsatile dynamics—sharp spikes rather than sinusoidal waves—reflect the nonlinear Hill kinetics and the cascade's rapid phosphorylation rates. Such oscillatory MAPK signaling has been observed in proliferation, differentiation, and stress responses, where periodic activation encodes distinct cellular decisions compared to sustained activation.

\subsection*{Adaptation: Perfect Responsiveness Despite Persistent Stimuli}

\textbf{The challenge:} Adaptation is perhaps the most paradoxical computational mode: the system must respond vigorously to stimulus onset yet return to baseline despite continued stimulus presence. This requires distinguishing \emph{change} from \emph{absolute level}. The circuit must amplify transient changes (high fold-change) while achieving near-perfect adaptation (returning close to baseline), yet these objectives are mathematically opposed—high amplification inherently conflicts with complete baseline return.

Our adaptive ERK cascade achieved 96.4\% adaptation with 3,219-fold amplification, operating at 98.8\% thermodynamic efficiency relative to the theoretical minimum steady-state elevation. When stimulated with a sustained 10-second growth factor pulse, ERK-PP spiked from 0.3 nM to 976 nM (peaking at 51 seconds), then adapted back toward baseline, reaching 36 nM by 180 seconds despite continued stimulus presence. This represents near-perfect change detection: the cascade reports \emph{when} the stimulus arrives but ignores its continued presence.

\begin{figure}[htbp]
\centering
\includegraphics[width=0.9\columnwidth]{adaptation_spike_timecourse.pdf}
\caption{Adaptive MAPK dynamics with near-perfect adaptation.}
\label{fig:adaptation}
\end{figure}

The adaptation timecourse (Figure~\ref{fig:adaptation}) shows the dual-pathway feedback architecture enabling both high amplification and efficient baseline return, operating at 98.8\% of the theoretical thermodynamic limit. The mechanism involves ERK-PP driving MKP synthesis (slow pathway providing persistent inhibition), while growth factor directly stimulates MKP production (fast pathway enabling anticipatory suppression). This architecture solves the temporal challenge of adapting to brief 10-second pulses—single-pathway feedback cannot accumulate inhibitor fast enough to prevent overshoot, but the direct feedforward enables rapid MKP accumulation synchronized with stimulus arrival.

\begin{figure}[htbp]
\centering
\includegraphics[width=0.9\columnwidth]{cascade_activation_timecourse.pdf}
\caption{Temporal activation sequence showing Raf→MEK→ERK propagation.}
\label{fig:cascade_activation}
\end{figure}

The cascade activation sequence (Figure~\ref{fig:cascade_activation}) demonstrates the temporal hierarchy of phosphorylation events over 0-40 seconds, with Raf activating first, followed by MEK, and finally ERK—illustrating how signals propagate through the three-tier architecture. The mathematical constraint analysis reveals why perfect adaptation with high amplification is impossible: achieving SRR $<$ 5× would require fold-change $<$ 115×, incompatible with MAPK cascade amplification capability. Our system operates at 98.8\% thermodynamic efficiency relative to this theoretical limit.

MKP inhibits all three phosphorylation steps (Raf, MEK, ERK) via Hill function:
\begin{equation}
\text{Inhibition} = \frac{K_m^n}{K_m^n + [\text{MKP}]^n}
\end{equation}
At the adapted state (MKP = 62 mM following the optimized 10s pulse), ERK-PP is suppressed to 0.036 mM (118× above baseline 0.0003 mM). The adaptation mechanism implements a type-1 incoherent feedforward loop (IFFL) \citep{mangan2003}: growth factor activates ERK directly (fast path: GF→Raf→MEK→ERK) and indirectly inhibits ERK via MKP synthesis (slow path: GF→MKP→ERK inhibition), creating temporal separation. The direct GF→MKP feedforward (rate=29) with fast Raf activation (rate=2.5) enables rapid coordinated accumulation during the brief pulse, preventing massive ERK-PP overshoot. Despite 3,219× amplification at peak, adaptation maintains numerical stability with an enhanced dephosphorylation rate of 42.0, operating at 98.8\% of the theoretical optimum efficiency. This efficiency metric reflects how close the system operates to the thermodynamic limit: given 96.4\% adaptation and 3,219× transient amplification, the theoretical minimum steady-state elevation is SRR$_{\text{min}}$=117× (below which adaptation would require thermodynamically infeasible parameters). Our system achieves SRR=118×, meaning 98.8\% efficiency—approaching the physical constraint where better adaptation would violate energy conservation.

Our computational results align quantitatively with experimental literature: bistability fold-change (577×) matches Ferrell's oocyte measurements (~500× separation), excitability threshold (15 nM) falls within Sasagawa's neuronal ERK activation range, oscillation frequency (20/min) reproduces Shankaran's EGF-stimulated pulsing, and adaptation efficiency (96.4\%) corresponds to Amit's observed homeostatic regulation. This quantitative agreement validates the model's biological fidelity while demonstrating the platform's capability to capture diverse nonlinear phenomena within a unified framework.

Each computational mode emerges from specific feedback architecture, quantified by two parameters: $\alpha$ (positive feedback strength: ERK-PP promotes PP2A degradation, accelerating its own activation) and $\beta$ (negative feedback strength: ERK-PP induces MKP synthesis, inhibiting cascade activity). The ratio $\alpha/\beta$ determines which computational mode the cascade exhibits. Table~\ref{tab:feedback_summary} summarizes feedback configurations across experiments:

\begin{table}[h]
\centering
\small
\caption{Feedback configurations and computational modes. $\alpha$ = positive feedback strength (ERK-PP→PP2A degradation), $\beta$ = negative feedback strength (ERK-PP→MKP synthesis). The ratio $\alpha/\beta$ acts as a tuning parameter determining cascade behavior.}
\label{tab:feedback_summary}
\begin{tabular}{@{}lrrr@{}}
\toprule
\textbf{Mode} & \textbf{$\alpha$} & \textbf{$\beta$} & \textbf{$\alpha/\beta$} \\
\midrule
Bistability & 15.0 & 1.0 & $>2$ \\
Excitability & 2.0 & 15.0 & 0.1 \\
Oscillations & $<$1 & 20.0 & $<$0.2 \\
Adaptation & 0 & 200 & -- \\
\bottomrule
\end{tabular}
\end{table}

These results demonstrate that the same molecular components—kinases, phosphatases, feedback loops—can be reconfigured to perform fundamentally different computational operations. The MAPK cascade is not a static relay, but a tunable information processor whose behavior emerges from the balance of feedback architectures.

\section*{Discussion}

These four computational modes—bistability, excitability, oscillations, and perfect adaptation—emerged from a single extended Petri net model through feedback modulation alone. Classical approaches require either phenomenological fits (losing mechanistic insight) or extensive ODE systems (struggling with information sensing). What enables this unified treatment?

Classical Petri nets (5-tuple: places, transitions, arcs, weights, marking) model local mass-transfer reactions but not non-local information sensing—kinases responding to extracellular growth factors require explicit arc chains, yet biological systems sense without consuming (quorum detection, ATP status, hormone gradients). Extending to 13-tuple formalism (adding signal places $\Psi$, rate functions $\Phi$, thermodynamic $\Delta$G, regulatory $\Sigma$) enables transitions to reference non-local places without arcs, revealing hierarchical information-flux layers invisible to classical formalism. In lambda phage, this formalism revealed 4 layers (environmental 1-8\% MI → stress RecA 43\% → integration CII 74\% → decision CI/Cro) where higher layers consume signal tokens from lower—cascade preemption, not mere integration. MAPK demonstrates hierarchical preemption through signal-directed computational mode selection: growth factor (signal place crossing membrane boundary) and MKP levels (signal place crossing cytoplasm/nucleus boundary) act as global information orchestrators that hierarchically determine which feedback architecture dominates—the signals collapse from global configuration (GF=high + MKP=low → bistability) to local deterministic cascade decisions (ERK-PP activation/deactivation). Different signal configurations preempt alternative modes: sustained high GF with rising MKP enforces adaptation, overriding bistability that would persist under high GF alone. While the implementations differ—lambda phage implements preemption through explicit token consumption (RecA depletes CII), whereas MAPK implements it through computational mode orchestration (GF/MKP levels select which feedback architecture dominates)—both exemplify information-flux hierarchy where higher-layer signals preempt lower-layer decisions. The unifying principle: signals at outer layers (environmental stress, extracellular growth factors) hierarchically determine which programs execute at inner layers (genetic decisions, cascade dynamics), creating multi-level control architectures. This is information-theoretic hierarchy: signals don't merely modulate—they orchestrate which computational program executes, with thermodynamic constraints enforcing trade-offs.

MAPK validates cross-kingdom generality of these principles. Feedback ratio $\alpha/\beta$ acts as a tuning parameter: $\alpha/\beta$>2 creates bistability (577× separation, Xenopus oocytes \citep{ferrell2009}), $\alpha/\beta$$\approx$0.1 produces excitability (15 nM threshold, 500× amplification), $\alpha/\beta$<0.2 enables oscillations (20 osc/min, proliferating cells \citep{shankaran2009}), and incoherent feedforward loop architecture achieves adaptation (96.4\%, 3,219× transient, 98.8\% thermodynamic efficiency: SRR$_{\text{min}}$=117×, achieved 118×, hormones \citep{amit2007}). Inherent trade-offs exist: bistability (persistent) conflicts with adaptation (baseline return), excitability (rapid reset) conflicts with oscillations (sustained)—these reflect mutually exclusive network architectures. Bistability requires coherent positive feedback ($\alpha$ > $\beta$, autocatalytic reinforcement), while adaptation requires incoherent feedforward loops ($\beta$ >> $\alpha$, fast direct + slow indirect paths with opposing signs). No single network topology satisfies both requirements simultaneously. Dual-pathway feedback (slow ERK-PP→MKP + fast GF→MKP) is necessary for adaptation—single-pathway feedback cannot accumulate inhibitor fast enough during 10s pulses.

Why use extended Petri nets versus ODEs? Three capabilities: (1) Non-local sensing ($\Psi$) captures information-theoretic signaling, (2) Thermodynamic validation ($\Delta$G) enforces energetic feasibility, (3) Weak independence ($\Sigma$) enables 80\% parallelizable transitions (versus 20\% in classical formalism). ODEs approximate biological systems through continuous mass-action kinetics—mathematically tractable but fundamentally treating molecular populations as continuous concentrations. Stochastic simulation engines (tau-leaping Skellam, Gillespie) model the physical reality: discrete molecular collisions, Brownian motion, individual particle trajectories. These are different physical modeling frameworks, not competing methods. Hybrid simulation combines stochastic sampling for low-copy regulators (where discrete events matter) with deterministic integration for abundant species (where continuous approximation is valid), capturing biological reality across concentration regimes.

\subsection*{Testable Predictions}

Our unified framework generates experimentally verifiable predictions: (1) Cells maintain specific $\alpha/\beta$ ratios characteristic of their computational mode—oocytes (bistable) should exhibit $\alpha/\beta$>2, proliferating cells (oscillatory) should show $\alpha/\beta$<0.2, and hormone-responsive cells (adaptive) should display $\beta$ >> $\alpha$ with dual-pathway architecture. These ratios should be measurable via quantitative Western blots tracking PP2A degradation rates ($\alpha$) and MKP synthesis rates ($\beta$) under different conditions. (2) Pharmacological perturbations should collapse computational modes as predicted: okadaic acid (PP2A inhibitor) should eliminate bistability by blocking positive feedback, forcing systems toward adaptation; MKP overexpression should shift oscillatory cells toward adaptation by increasing $\beta$; and MKP knockdown should convert adaptive cells toward bistability/excitability. (3) The thermodynamic efficiency constraint predicts adaptation systems operate near SRR$_{\text{min}}$ limits—measuring steady-state ERK-PP elevation and transient amplification across cell types should reveal consistent proximity to theoretical minima, with deviations indicating non-optimal parameter regimes or alternative network architectures. (4) Excitability thresholds should correlate with ATP availability: lowering [ATP] should increase the growth factor threshold for all-or-nothing responses, as phosphorylation becomes energy-limited. These predictions distinguish our framework from phenomenological models lacking thermodynamic constraints.

\subsection*{Thermodynamic Validation Framework}

Thermodynamic validation operates through compound mappings: place names (ATP, ADP, Raf) are mapped to KEGG compound IDs (C00002, C00008) enabling Gibbs free energy calculations from biochemical databases. For each reversible reaction, the validator computes equilibrium constant K$_{eq}$ = exp($-\Delta$G°/RT) from standard formation energies, then checks whether kinetic rate ratios (k$_{forward}$/k$_{reverse}$) match K$_{eq}$ within tolerance. Reactions with inconsistent ratios are flagged as thermodynamically infeasible—they would violate energy conservation under cellular conditions ([ATP]=5 mM, [ADP]=1 mM, pH 7.0, 298 K). This constraint rejected parameter sets where phosphorylation proceeded with $\Delta$G > 0 (endergonic without ATP coupling) or where k$_{forward}$/k$_{reverse}$ deviated >10× from K$_{eq}$ (kinetically implausible). The 98.8\% thermodynamic efficiency for adaptation refers to steady-state response ratio SRR=118× approaching the theoretical minimum SRR$_{min}$=117× imposed by 96.4\% adaptation with 3,219× transient amplification—cells operate near physical efficiency limits where higher fold-change or better adaptation would require thermodynamically infeasible parameter regimes.

\section*{Conclusion}

Extending classical Petri nets (5-tuple mass-transfer) to 13-tuple Bio-PN (adding signal places $\Psi$, thermodynamics $\Delta$G, regulatory $\Sigma$) transforms notation into a discovery engine: it predicts hierarchical cascade preemption through information-flux layers. Lambda phage: 4 layers where RecA stress (43\% MI) consumes developmental CII signals (74\% MI) with 2.01× priority (UV→71\% lytic). MAPK validates cross-kingdom generality: a single Raf→MEK→ERK cascade exhibits bistability (577× separation, $\alpha/\beta$>2), excitability (500× amplification, $\alpha/\beta$$\approx$0.1), oscillations (20/min, $\alpha/\beta$<0.2), and adaptation (96.4\%, 3,219× transient, 98.8\% thermodynamic efficiency SRR$_{\text{min}}$=117×) through feedback modulation. Trade-offs are thermodynamic: bistability conflicts with adaptation, excitability with oscillations—incompatible architectures, not tuning limits.

Why extended Petri nets? Non-local sensing ($\Psi$) captures information-theoretic signaling without arc connections, thermodynamic validation ($\Delta$G) enforces energy constraints, and weak independence ($\Sigma$) enables 80\% parallelism (versus 20\% in classical formalism). Dual-pathway adaptation (slow ERK-PP→MKP + fast GF→MKP) operates at SRR=118× (theoretical minimum 117×)—cells maintain homeostasis at physical efficiency limits.

The discovery-to-validation arc establishes the principle: biological signaling implements hierarchical control through signal places—global information orchestrators crossing compartment boundaries that hierarchically determine local deterministic decisions. In lambda phage, RecA stress signals consume CII developmental tokens (explicit token depletion). In MAPK, growth factor and MKP signals orchestrate which feedback architecture dominates (computational mode selection). Both mechanisms manifest information-flux hierarchy: signals at higher layers (environmental, membrane-crossing) preempt decisions at lower layers (genetic, cascade-level), collapsing from global configuration to local execution. From viral lysogeny (lambda phage) to bacterial sporulation (\emph{Bacillus subtilis}) to eukaryotic signaling (MAPK), Signal Hierarchical Petri Nets capture this information-theoretic control invisible to classical mass-action formalism—cells orchestrate survival through hierarchical signal preemption, now tractable through thermodynamically-grounded modeling platforms.

\section*{Methods}

\subsection*{Signal Hierarchical Petri Net Formalism: A Theoretical Journey}

Classical Petri nets (5-tuple $\text{PN} = (P, T, F, W, M_0)$: places/tokens, transitions/reactions, arcs/flow, weights/stoichiometry, marking/distribution) excel at mass-transfer modeling through bipartite token-flow graphs (state equation $M' = M + C \cdot \sigma$, incidence matrix $C = F^+ - F^-$), which have founded biological modeling \citep{chaouiya2007, heiner2008}. Limitation: transitions sense only arc-connected places—growth factors informing kinases require explicit arc chains, yet biological sensing is non-local (quorum detection, ATP status). Classical Petri nets force artificial modeling: either create explicit arcs (geometrically complex) or ignore signals (losing realism).

Theoretical breakthrough: extend transitions from simple firing to rate functions $\Phi: T \rightarrow \mathbb{R}^+$ referencing any place. Signal places $\Psi$ are formalized via set subtraction: $\Psi(t) = \text{Referenced}(\Phi(t)) \setminus ({}^\bullet t \cup t^\bullet \cup \Sigma(t))$—places in the rate function but not arc-connected (consumed ${}^\bullet t$, produced $t^\bullet$, regulatory $\Sigma(t)$) are sensed. Automatic detection: parse rates, subtract arc-connected places, and designate the remainder as signals. Concrete example from MAPK: the Raf activation transition has rate function $\Phi(t_{\text{Raf-activation}}) = k_{\text{basal}} + k_{\text{GF}} \cdot [\text{GF}] \cdot [\text{Raf}] / (K_m + [\text{Raf}])$, where GF (growth factor) appears in the rate but is not consumed ($\text{GF} \notin {}^\bullet t_{\text{Raf-activation}}$), not produced ($\text{GF} \notin t_{\text{Raf-activation}}^\bullet$), and not connected via regulatory arcs ($\text{GF} \notin \Sigma(t_{\text{Raf-activation}})$). Therefore, GF is automatically classified as a signal place: $\text{GF} \in \Psi(t_{\text{Raf-activation}})$. This enables non-local sensing where membrane-bound growth factors inform cytoplasmic kinase activation without requiring explicit arc chains through receptors and transport machinery.

13-tuple Bio-PN $(P, T, F, W, M_0, \Psi, \Phi, \Delta G, \Sigma, \text{ATP}, \text{ADP}, K, \tau)$ integrates: classical 5-tuple (mass foundation), signal places $\Psi$ (non-local sensing), rate functions $\Phi$ (Michaelis-Menten/Hill/mass-action), Gibbs energy $\Delta G$ (thermodynamic feasibility, ATP-coupled $\Delta G < 0$), regulatory arcs $\Sigma$ (test/inhibitor enabling weak independence: competitive 20\% sequential, regulatory 80\% parallel), ATP/ADP (energy currency), constants $K$ (kinetics), delays $\tau$ (transcription).

Discovery validation: lambda phage signal place detection revealed a 4-layer hierarchy (environmental → RecA stress 43\% MI → CII integration 74\% → CI/Cro decision), where higher layers consume signal tokens (cascade preemption), thermodynamically validated ($\Delta$G determines activation), and experimentally confirmed (UV→71\% lytic, 2.01× priority). MAPK instantiates the framework: spatial compartments (membrane/cytoplasm/nucleus), hybrid simulation (tau-leaping Skellam for reversible reactions + continuous deterministic for abundant species + immediate/timed for regulatory events), thermodynamic $\Delta$G validation, and weak independence (80\% parallelizable). This captures biological reality: discrete events, spatial organization, energy constraints, and information sensing.

\subsection*{Model Architecture}

The MAPK cascade model comprises three tiers: Raf (MAPKKK), MEK (MAPKK), and ERK (MAPK), each existing in unphosphorylated and phosphorylated states. Growth factor (GF) activates Raf at the membrane. Active Raf-P phosphorylates MEK to MEK-P, which phosphorylates ERK to ERK-P and ERK-PP (doubly phosphorylated). Phosphatases PP2A (acting on Raf and MEK) and MKP (acting on ERK) catalyze dephosphorylation.

\begin{figure}[htbp]
\centering
\includegraphics[width=0.95\columnwidth]{erk_cascade_adaptation.pdf}
\caption{MAPK cascade architecture with feedback loops.}
\label{fig:cascade_architecture}
\end{figure}

The cascade architecture (Figure~\ref{fig:cascade_architecture}) shows how ERK-PP drives MKP synthesis (slow negative feedback) while GF directly stimulates MKP production (fast feedforward). MKP globally inhibits all three phosphorylation tiers via Hill function, enabling near-perfect adaptation with high transient amplification. Feedback architectures were implemented as follows:
\begin{itemize}
\item \textbf{Positive feedback:} ERK-PP promotes PP2A degradation with coefficient $\alpha$, creating autocatalytic activation
\item \textbf{Negative feedback:} ERK-PP induces MKP synthesis with coefficient $\beta$ and Hill kinetics ($n=2$, $K_d=10$ nM)
\item \textbf{Dual-pathway adaptation:} Combined slow ERK-PP→MKP feedback ($\beta=200$) with fast direct GF→MKP feedforward (rate 29)
\end{itemize}

MKP inhibits all three tiers via Hill function: $\text{Inhibition} = K_m^n / (K_m^n + [\text{MKP}]^n)$ with $n=3$ and $K_m=75$ $\mu$M.

\subsection*{Simulation Protocol}

Hybrid stochastic-continuous simulation used SHYPN's tau-leaping algorithm with Skellam sampling for reversible reactions (synthesis/degradation), combined with continuous mass-action kinetics for abundant species (phosphorylation cascades) and immediate transitions for instantaneous regulatory events. Time step $\Delta t = 0.0001$ seconds balanced stochastic accuracy (tau-leaping condition: expected firings $<10$ per step) with computational efficiency. Simulations ran for 180-300 seconds depending on the phenomenon (bistability: 300s, excitability: 180s, oscillations: 180s, adaptation: 180s). Stochastic simulation models discrete molecular events—individual phosphorylation reactions occur through random collisions governed by Brownian motion, with reaction probabilities computed from molecular concentrations and kinetic constants. For low-copy species (e.g., active phosphatases, regulatory signals), discrete stochasticity dominates: 10 molecules behave fundamentally differently than a 10 nM concentration field. For high-abundance species (cascade intermediates at $\mu$M concentrations), the law of large numbers makes continuous approximation valid. This hybrid approach captures the physical reality: biology operates across concentration regimes where both discrete events and continuous flows matter. Thermodynamic validation computed $\Delta G$ for all ATP-coupled reactions using standard conditions (298 K, pH 7.0, [ATP]=5 mM, [ADP]=1 mM).

\subsection*{Analysis Methods}

\textbf{Bistability:} Basin of attraction computed by simulating from 100 initial conditions spanning ERK-PP = 0-600 nM, classifying final states as HIGH (>300 nM) or LOW (<100 nM). Hysteresis verified by applying stimulus pulses of varying duration and measuring persistence.

\textbf{Excitability:} Dose-response curves measured peak ERK-PP amplitude as function of GF concentration (0-100 nM). Threshold defined as GF concentration producing half-maximal activation. Refractoriness assessed via paired pulses with varying interpulse intervals.

\textbf{Oscillations:} Fourier analysis computed dominant frequency peaks using FFT. Period measured as mean time between successive ERK-PP peaks detected using peak-finding algorithm with prominence threshold 100 nM. Period stability quantified as coefficient of variation across all cycles.

\textbf{Adaptation:} Metrics included: (1) Adaptation percentage $A = 100 \times ([\text{ERK-PP}]_{\text{peak}} - [\text{ERK-PP}]_{\text{adapted}}) / ([\text{ERK-PP}]_{\text{peak}} - [\text{ERK-PP}]_{\text{baseline}})$; (2) Fold-change $F = [\text{ERK-PP}]_{\text{peak}} / [\text{ERK-PP}]_{\text{baseline}}$; (3) Steady-state response ratio $\text{SRR} = [\text{ERK-PP}]_{\text{adapted}} / [\text{ERK-PP}]_{\text{baseline}}$; (4) Adaptation timescale $\tau_{\text{adapt}}$ as the time to reach 10\% above baseline post-peak. The theoretical minimum SRR was computed as $\text{SRR}_{\text{min}} = 1 + (1 - A/100) \times (F - 1)$.

\subsection*{Model Parameters}

\textbf{Bistability:} Positive feedback $\alpha = 15.0$ (ERK-PP promotes PP2A degradation), negative feedback $\beta = 5.0$, basal PP2A 20 nM, GF pulse 100 nM for 30s.

\textbf{Excitability:} Balanced feedback $\alpha = 2.0$, $\beta = 15.0$, MKP synthesis Hill $n=3$ with $K_m = 75$ $\mu$M, MKP degradation 0.05.

\textbf{Oscillations:} ERK-PP→MKP synthesis 20.0 with Hill kinetics ($n=2$, $K_d=10$ nM), ERK dephosphorylation 8.0, MKP degradation 0.02, sustained GF 50 nM.

\textbf{Adaptation:} Dual-pathway with slow ERK-PP→MKP synthesis (200.0) and fast GF→MKP feedforward (29.0). Raf activation 2.5, cascade phosphorylation 0.3, ERK dephosphorylation 42.0, MKP global inhibition Hill $n=3$ with $K_m=75$ $\mu$M, GF pulse 0.1 mM for 10s.

\textbf{Common:} Phosphorylation consumed ATP ($K_m = 100$ $\mu$M). Basal rates: Raf activation 1.0, MEK/ERK phosphorylation 0.3, dephosphorylation 1.0. Initial conditions: all species unphosphorylated, [ATP]=5 mM, [ADP]=1 mM.

\section*{Funding}

This research received no specific grant from any funding agency in the public, commercial, or not-for-profit sectors.

\section*{Acknowledgments}

The author thanks the Universidade Federal de Santa Catarina for computational resources and the systems biology community for foundational work on Petri net formalisms.

\section*{Author Contributions}

E.S. conceived the study, developed the Signal Hierarchical Petri Net formalism, implemented the computational models, performed all simulations, analyzed the data, and wrote the manuscript.

\section*{Competing Interests}

The author declares no competing interests.

\section*{Data Availability}

Model files (.shy), simulation data, analysis scripts, and figure generation code are available at: \url{https://github.com/simao-eugenio/shypn}.

\section*{Software Availability}

Signal Hierarchical Petri Net models and simulations were performed using SHYPN 2.0 (Signal Hierarchical Petri Nets platform). The software implements the 13-tuple Bio-PN formalism with hybrid stochastic-deterministic simulation, thermodynamic validation, and spatial organization. Source code, documentation, and installation instructions are freely available under open-source license at \url{https://github.com/simao-eugenio/shypn}.

\bibliographystyle{plainnat}
\begin{thebibliography}{99}
\setlength{\itemsep}{3pt}

\bibitem{chaouiya2007}
Chaouiya, C. (2007).
\newblock Petri net modelling of biological networks.
\newblock \textit{Brief. Bioinform.} 8(4):210--219.

\bibitem{heiner2008}
Heiner, M., Gilbert, D., and Donaldson, R. (2008).
\newblock Petri Nets for Systems and Synthetic Biology.
\newblock \textit{Lect. Notes Comput. Sci.} 5016:215--264.

\bibitem{mangan2003}
Mangan, S. and Alon, U. (2003).
\newblock Structure and function of the feed-forward loop network motif.
\newblock \textit{Proc. Natl. Acad. Sci. USA} 100(21):11980--11985.

\bibitem{simao2025thermodynamic}
Simão, E. (2025).
\newblock Thermodynamic Constraints Drive Hierarchical Preemption in Cellular Decision-Making: A Hybrid Petri Net Framework with Application to \textit{Bacillus subtilis} Sporulation.
\newblock arXiv:2601.04335.

% MAPK Historical Foundation
\bibitem{marshall1994}
Marshall, C.J. (1994).
\newblock MAP kinase kinase kinase, MAP kinase kinase and MAP kinase.
\newblock \textit{Curr. Opin. Genet. Dev.} 4(1):82--89.

% MAPK Bistability
\bibitem{xiong2003}
Xiong, W. and Ferrell, J.E. (2003).
\newblock A positive-feedback-based bistable 'memory module' that governs a cell fate decision.
\newblock \textit{Nature} 426(6965):460--465.

\bibitem{ferrell2009}
Ferrell, J.E. (2009).
\newblock Signaling motifs and Weber's law.
\newblock \textit{Mol. Cell} 36(5):724--727.

% MAPK Oscillations
\bibitem{kholodenko2000}
Kholodenko, B.N. (2000).
\newblock Negative feedback and ultrasensitivity can bring about oscillations in the mitogen-activated protein kinase cascades.
\newblock \textit{Eur. J. Biochem.} 267(6):1583--1588.

\bibitem{shankaran2009}
Shankaran, H., Ippolito, D.L., Chrisler, W.B., Resat, H., Bollinger, N., Opresko, L.K., and Wiley, H.S. (2009).
\newblock Rapid and sustained nuclear-cytoplasmic ERK oscillations induced by epidermal growth factor.
\newblock \textit{Mol. Syst. Biol.} 5:332.

\bibitem{albeck2013}
Albeck, J.G., Mills, G.B., and Brugge, J.S. (2013).
\newblock Frequency-modulated pulses of ERK activity transmit quantitative proliferation signals.
\newblock \textit{Mol. Cell} 49(2):249--261.

% MAPK Excitability
\bibitem{sasagawa2005}
Sasagawa, S., Ozaki, Y., Fujita, K., and Kuroda, S. (2005).
\newblock Prediction and validation of the distinct dynamics of transient and sustained ERK activation.
\newblock \textit{Nat. Cell Biol.} 7(4):365--373.

% MAPK Adaptation
\bibitem{amit2007}
Amit, I., Citri, A., Shay, T., Lu, Y., Katz, M., Zhang, F., Tarcic, G., Siwak, D., Lahad, J., Jacob-Hirsch, J., Amariglio, N., Vaisman, N., Segal, E., Rechavi, G., Alon, U., Mills, G.B., Domany, E., and Yarden, Y. (2007).
\newblock A module of negative feedback regulators defines growth factor signaling.
\newblock \textit{Nat. Genet.} 39(4):503--512.

\bibitem{mackeigan2005}
Mackeigan, J.P., Murphy, L.O., and Blenis, J. (2005).
\newblock Sensitized RNAi screen of human kinases and phosphatases identifies new regulators of apoptosis and chemoresistance.
\newblock \textit{Nat. Cell Biol.} 7(6):591--600.

\bibitem{qiao2007}
Qiao, L., Nachbar, R.B., Kevrekidis, I.G., and Shvartsman, S.Y. (2007).
\newblock Bistability and oscillations in the Huang-Ferrell model of MAPK signaling.
\newblock \textit{PLoS Comput. Biol.} 3(9):e184.

% Recent MAPK Reviews
\bibitem{cargnello2011}
Cargnello, M. and Roux, P.P. (2011).
\newblock Activation and function of the MAPKs and their substrates, the MAPK-activated protein kinases.
\newblock \textit{Microbiol. Mol. Biol. Rev.} 75(1):50--83.

\bibitem{lavoie2020}
Lavoie, H., Gagnon, J., and Therrien, M. (2020).
\newblock ERK signalling: a master regulator of cell behaviour, life and fate.
\newblock \textit{Nat. Rev. Mol. Cell Biol.} 21(10):607--632.

\bibitem{lake2016}
Lake, D., Corrêa, S.A., and Müller, J. (2016).
\newblock Negative feedback regulation of the ERK1/2 MAPK pathway.
\newblock \textit{Cell. Mol. Life Sci.} 73(23):4397--4413.

\end{thebibliography}

\end{document}
